\chapter{給第一本小說的五個建議}

這場講座並不是黑板式的數學講課，但它以另一種方式極其精確。它從一個數值限制開始：在一個月內寫出五萬字；從這個限制出發，Sanderson 逐步建起一連串實用構造：借來一套結構、打開一條通向角色的管道、用欲望與需求製造壓力、讓整個故事圍繞「進展」來組織，最後則在真正坐下寫之前，先把心智調到準備就緒的狀態。這個次序很重要，因為每一步都在回應前一步仍然不夠的地方。

\begin{remark}
本章沒有任何經過驗證的講座截圖留存。因此，下文中的公式與唯一的一張示意圖，都只能依據逐字稿重建，而不是從可見的黑板內容轉錄而來。
\end{remark}

\section{挑戰、授權與五個捷徑的設置}

我們就從 Sanderson 開始的地方開始：把 National Novel Writing Month 當成一個實作挑戰，而不是一條文學教義。這個挑戰的工作目標，是在一個月內寫出一部小說，而這裡對「小說」的定義也被放得足夠寬，好讓挑戰本身成立。Sanderson 立刻補充說，五萬字只是一個粗略的操作數字，並不是對形式的理論；他自己的許多小說都遠比這長。但這項挑戰依然是嚴肅的，因為它同時既做得到，又不容易。

\begin{align}
\text{NaNoWriMo 目標} &= 50{,}000\ \text{字},\\
\text{每日進度} &\approx \frac{50{,}000}{30} \approx 1{,}667 \approx 1{,}700\ \text{字／日}.
\end{align}

這項練習的目的也同樣重要。我們是要藉它走出卡住的狀態、關掉內在編輯器，然後開始寫。Sanderson 又用個人例子加強這一點：在出版之前，他做過這個挑戰很多次，甚至還把 \textit{The Way of Kings} 最初的書寫追溯到某一次這樣的月度衝刺。因此，講座一開始不是從抽象講起，而是從一個可工作的限制與一則能證明此限制確有生產力的見證講起。

但同樣重要的是，他幾乎立刻給出了一道許可：這件事未必適合每一位作者。這個挑戰可能與某個人的寫作心理相衝突，也可能讓最後寫出來的頁面，比起慢慢寫出的頁面更差。這種保留，決定了整場講座的語氣。這些是方法，而不是戒律。

\subsection{Question \& Answer}

\paragraph{Question.} 這個挑戰什麼時候有用？又什麼時候該放棄？

\paragraph{Answer.} 當一個外在而明確的限制能幫助我們打破遲疑、壓低內在編輯器的聲音，並累積頁面時，它就是有用的；而當這個限制明顯開始妨礙作品、違逆我們的寫作心理，或讓文字變得更差而不是更自由時，我們就該放棄它。

只有在把這片地清出來之後，Sanderson 才提出這場講座真正的承諾：如果你確實想試試看，那麼他會給你五個 hacks，五種即使準備不多，也能開始一部小說的實際方法。

\section{借用結構，作為輔助輪}

第一個捷徑，從一種作者常有的衝動開始。我們看了一部電影、讀了一個故事、聽說了一件事，然後心想：我能不能也寫出類似的東西？Sanderson 的回答不是去複製那個故事，而是分析：在那類故事裡，究竟是什麼結構在運作。這種情況不會發生在每一本小說上，也不會發生在每一位作者身上；但一旦發生，它就能立刻給我們抓地力。

\begin{definition}
所謂借來的結構，是指從某個故事或某個子類型中抽取出的功能性模式，然後用不同的角色、不同的利害關係與不同的弧線，重新將其建造出來。
\end{definition}

講座裡最清楚的例子，是劫案故事。吸引力並不只是它的犯罪外殼，而在於它的結構：一個困難的問題，一組彼此互補的專家角色，準備工作，以及一個必須讓這些專長最終協同運作的解決方式。

\begin{equation}
\text{受喜愛的故事} \to \text{標誌性場景} \to \text{基本結構} \to \text{新故事}.
\end{equation}

在 Sanderson 的劫案例子裡，這套結構可以示意地壓縮成

\begin{equation}
\text{問題} \to \text{招募} \to \text{準備} \to \text{協作式解決}.
\end{equation}

\paragraph{A worked extraction rule.} 講座的方法可以寫成：
\begin{enumerate}
\item 從一個我們喜歡的故事或子類型出發。
\item 問：這類故事之所以令人滿足，究竟是因為什麼？
\item 抽出其中那些標誌性的節拍或場景。
\item 把這些節拍煮沸成一個更基本的結構。
\item 再用新的角色、新的問題與新的角色弧線，把這個結構重新搭起來。
\end{enumerate}

一旦我們看見這一點，結構就開始從原本的例子身上鬆開。Sanderson 接著又用移植來加強這個觀念：我們完全可以把自己在 Regency romance 裡喜歡的一種結構，搬進一部 Western 裡。這樣的移植保留了骨架，卻換掉了血肉，也能防止借用凝固成模仿。

\subsection{Question \& Answer}

\paragraph{Question.} 我們怎樣借用一套結構，而不反過來被它奴役？

\paragraph{Answer.} 借的是模式，而不是身分。原來那個故事提供的是一條承重序列，但它的表面世界、角色、核心問題與弧線，都必須被替換掉。借來的結構像輔助輪：它確實支撐你，但不是一座永遠的牢籠。

這最後一句非常關鍵。Sanderson 說，在我們想快速寫作時，這種方法可能幫上大忙；但他同時也堅持，即使是專業作家也會這樣借結構。因此，這不只是補救性的工具，而是一種正當的方式，用來讓想像穩定下來。

\section{從獨白開始}

一旦外部支架就位，講座便從結構收窄到進入點。要怎樣足夠快地進入一個角色，快到能夠開始寫？Sanderson 的第二個答案，是獨白。即使最後完成的小說不會採用第一人稱形式，我們仍然可以讓角色像是在直接對我們說話一樣，去解釋生命中某件重要的事。

\begin{equation}
\text{獨白} \to \text{聲音} + \text{歷史} + \text{可用碎片}.
\end{equation}

這裡的重點是發現，而不是最終形式。Sanderson 舉例說，我們可以要求一個角色：如果要你用五頁紙解釋人生，你會怎麼說？這五頁未必會真正印進最後成書的版本裡，但它們會暴露出措辭、重點、記憶、自我欺瞞與態度。從這個意義上說，這段獨白首先不是寫給讀者，而是寫給作者。

接著，講座又把這個練習往前推一步。用這種方式產出的材料，可以變成題記、日記碎片，或章首其他細小的前置文本。如果這種形式本身特別有生產力，它甚至還可能提示出一整部書信體小說。

\begin{definition}
書信體小說，是由作品內部文件——例如日記、信件，或其他類似文書——組成的小說。
\end{definition}

Sanderson 那句隨口的描述在這裡很有幫助：從某種意義上說，這就像小說版的 found footage。這個類比是鬆的，但它能夠很快抓住形式上的關鍵。

\subsection{Question \& Answer}

\paragraph{Question.} 即使最後完成的小說不是第一人稱，第一人稱獨白也有用嗎？

\paragraph{Answer.} 有。在這裡，獨白不是一項關於視角的教條選擇，而是一種發現工具。我們用它來聽見角色、抽出聲音與歷史，然後才決定：其中究竟有什麼會存活進最後的書裡。

這就是講座的第二個主要模式：一個暫時性的練習，可以始終保持暫時性；但如果它所產出的聲音太有價值，不捨丟棄，它也可能直接轉化為正式材料。

\section{欲望、需求、障礙，以及選對視角角色}

第三個捷徑，從聲音轉向壓力。Sanderson 問了四個彼此相扣的問題：角色想要什麼？角色需要什麼？這兩者有何不同？以及，為什麼角色不能輕易得到任何一個？到了這裡，講座的結構變得更乾淨了。情節不再是從角色外部施加進去的東西，而是從有意識的追求與更深層必要之間的不吻合處生成出來的。

讓我們先替這裡的核心量命名：
\begin{align}
W &= \text{角色想要的東西},\\
N &= \text{角色需要的東西},\\
O &= \text{阻擋兩者的障礙}.
\end{align}

其中最關鍵的區分，是前兩者的不相等：

\begin{equation}
W \neq N.
\end{equation}

一旦這條裂縫被打開，情節壓力便會隨之而來：

\begin{equation}
\text{情節壓力} = \text{對 } W \text{ 與 } N \text{ 的障礙}.
\end{equation}

\paragraph{A worked update rule.} Sanderson 從角色生成情節的程序，可以壓縮成：
\begin{enumerate}
\item 問角色想要什麼。
\item 問角色需要什麼。
\item 問這兩者有何不同。
\item 問角色為什麼不能乾脆得到其中任一項。
\item 把這些答案轉化為障礙。
\item 讓這些障礙決定情節最初的運動方式。
\end{enumerate}

這也就是為什麼講座會堅持：故事應當以角色為中心。若情節的轉折來自 \(W\)、\(N\) 與 \(O\) 周圍的壓力，那麼角色就會被安置在小說機械的內部，而不是站在它的外面。

接著，Sanderson 診斷出一個初學者常犯的錯誤。我們可能選錯了視角角色：某個人只是站在旁邊看事件發生，而真正的戲劇其實發生在另一個人身上。這樣一來，名義上的主角就成了別人的故事的外部觀察者。

\subsection{Question \& Answer}

\paragraph{Question.} 究竟哪個角色才真正應該承載故事？

\paragraph{Answer.} 不是那個被動旁觀的人。故事應當屬於改變最多、承受最多衝突，或最積極去追求其所欲之物的那個角色。若真正的壓力落在另一個人物身上，那麼那個人其實更接近真正的主角。

講座把這個診斷當成一座橋。Sanderson 想像作者接著會說：好，現在我知道角色想要什麼，也知道他需要什麼，但我還是沒有結構。這個抱怨，便打開了整場談話中最有形式感的一段。

\section{承諾、進展、回報，以及進展的標示方式}

第四個捷徑，明白地被提出來作為對「我還是沒有結構」的回答。Sanderson 說，他關於這個主題其實有整整一場課；但如果把它煮濃，故事是由三件事構成的：

\begin{equation}
\text{故事} = \text{承諾} \to \text{進展} \to \text{回報}.
\end{equation}

這個三元結構值得被放在本章的核心附近。承諾，是故事說它將提供哪一類運動與滿足；回報，則是那份承諾的兌現。但夾在中間的那個巨大項——進展——才構成故事的大部分。讀者之所以持續翻頁，是因為他們能感到自己正朝向某樣東西移動。

在這裡，Sanderson 補上一個出人意料地有用的觀察：就連書本身的物理形狀，也會提供一種粗略的進展感。隨著頁數翻動，我們會離結尾更近。但這種粗糙的運動仍然不夠。情節必須加以補強，給我們一條更具體的前進通道。

\begin{figure}
\centering
\begin{tikzpicture}
\draw (-5.2,0.8) rectangle (-2.4,-0.8);
\node at (-3.8,0.12) {承諾};

\draw (-1.0,0.8) rectangle (1.8,-0.8);
\node at (0.4,0.12) {進展};

\draw (3.2,0.8) rectangle (6.0,-0.8);
\node at (4.6,0.12) {回報};

\draw[->] (-2.4,0) -- (-1.0,0);
\draw[->] (1.8,0) -- (3.2,0);

\draw (-2.2,-3.0) rectangle (0.6,-1.6);
\node at (-0.8,-2.0) {旅行};
\node at (-0.8,-2.42) {進展};

\draw (1.2,-3.0) rectangle (4.4,-1.6);
\node at (2.8,-2.0) {資訊};
\node at (2.8,-2.42) {進展};

\draw[->] (0.4,-0.8) -- (0.4,-1.2);
\draw[->] (0.4,-1.2) -- (-0.8,-1.6);
\draw[->] (0.4,-1.2) -- (2.8,-1.6);
\end{tikzpicture}
\end{figure}

這是一張完全基於逐字稿的示意圖。它並不聲稱自己有講座畫面作為視覺背書。它只是把 Sanderson 所做的區分顯示出來：故事的一般結構是一回事，而「進展」則可以沿著不同通道被測量。

他的第一個例子是地理性的。在一場追尋中，我們可以用空間來衡量運動：

\begin{equation}
\text{旅行進展} : \text{起點} \to \text{目的地}.
\end{equation}

講座很小心，並不把進展等同於一路順利往前。角色可能被岔開、被迫後退，或被阻擋在自己原本以為的路線之外。但這些逆轉之所以仍然可讀，正是因為目的地仍然夠穩定，讓我們知道「更接近」究竟意味著什麼。

第二個例子則是資訊性的。在推理故事裡，進展主要不是空間上的，而是認識論上的：

\begin{equation}
\text{資訊進展} : \text{線索稀少} \to \text{線索更多} \to \text{圖像更清晰}.
\end{equation}

因此，假線索與誤導並不會摧毀結構；只要讀者仍然能感到整體圖像正在變清楚，它們反而屬於結構的一部分。

從這裡，Sanderson 走向他對節奏的診斷。一個故事之所以顯得慢，未必是因為什麼事都沒發生，而是因為它所承諾的那種「進展」沒有被充分標示出來。

\begin{equation}
\text{被清楚標示的進展} \Rightarrow \text{翻頁的動量}.
\end{equation}

而反方向的情況，也同樣重要：

\begin{equation}
\text{錯位的標示} \Rightarrow \text{被感知為節奏不佳}.
\end{equation}

\paragraph{Transcript note.} 逐字稿在這裡的例子相當粗糙，但想表達的結構點非常清楚。假如一本書最強烈地被標示成一場冒險，而作者自己私下其實把它想成一部愛情小說，那麼讀者就會在冒險軸上衡量進展；如果那條被承諾的通道沒有被持續推進，讀者就會感到停滯。

\subsection{Question \& Answer}

\paragraph{Question.} 明明場景裡一直有事發生，為什麼一個故事仍然會顯得慢？

\paragraph{Answer.} 因為事件本身並不自動產生動量。動量需要一條可見的進展通道。若我們承諾的是旅程，讀者就必須感到空間上的接近；若我們承諾的是推理，讀者就必須感到資訊在累積。當這條通道太弱，或者它與故事最強的標示彼此錯位時，整本書就會顯得靜止，即使場景其實很活躍。

Sanderson 在本節最後的建議，既實際又精確：把那條較大的進展通道再拆成較小而可辨識的單位，像灑下麵包屑那樣，讓讀者不只在整體上能感到前進，也能在局部上感到前進。

\section{為心智做預熱，先於頁面開始書寫}

第五個也是最後一個捷徑，把焦點從故事機械轉回工作機械。若我們的目標是一天大約 \(1{,}700\) 字，那就不能每一次都在完全冰冷的狀態下開工。Sanderson 的解法，是在真正坐下來寫之前，先替心智做預熱。

\begin{equation}
\text{頁外排演} \to \text{頁上執行}.
\end{equation}

講座對這件事十分具體。那些可利用的時間可能是通勤、運動、洗碗、割草，或任何讓身體忙碌、卻讓心智相對自由的活動。我們不必讓這些時間散掉，而可以用它們來問：下一場景究竟該完成什麼？是什麼能讓它不只是例行公事，而變得鮮明？又是什麼會讓它成為整本書裡值得記住的一場？

\paragraph{A worked routine.} Sanderson 的方法可以概括為：
\begin{enumerate}
\item 先決定下一個場景是什麼。
\item 問這個場景裡必須發生什麼。
\item 問這一場究竟有什麼值得寫，而不只是履行程序。
\item 在心裡先排演它好幾次。
\item 等真正坐下來時，腦中已經先知道下一段運動會怎麼展開。
\end{enumerate}

他甚至還補上一種情緒上的微調：音樂也許能幫助想像進入場景本身所要求的情緒頻率。這裡的重點不是裝飾，而是減少啟動摩擦。到我們真正坐到桌前時，那個場景其實已經在心裡開始了。

最後，講座又把這個建議推廣到「一個月小說挑戰」之外。即使日常生活只允許我們在週末寫作，一週中的頁外預熱仍然能夠替週末的幾個小時做準備。因此，整場講座的最後動作在結構上也很工整：在教完我們如何讓故事前進之後，Sanderson 最後教的是如何讓作者自己也能一直前進。

\section{Summary}

這場講座比它輕鬆的語調更有形式紀律。我們先從一個數字目標與一種許可結構開始：可以快速寫、可以關掉內在編輯器，但也不該把挑戰本身當成神明，只要它開始反過來傷害你的寫作心理。接著，我們借來外部結構，透過獨白發現角色聲音，從欲望與需求的裂縫中生成情節，再把整個故事壓縮成承諾、進展與回報，最後則以「先於頁面思考」的方式，支撐起這整套工程。

因此，本章最核心的教訓其實是程序性的。當我們在壓力之下開始一部小說時，只要給想像力足夠的穩定結構——一個外在限制、一個可回收的模式、一個角色引擎、一條可見的進展通道，以及一種能夠減少意圖與執行之間落差的日常實踐——小說就可以啟動。Sanderson 就是以這樣的順序帶我們走完講座，而也正是這個順序，讓這份筆記讀起來像講座本身在展開，而不是幾句彼此斷裂的格言。